% SPDX-License-Identifier: Apache-2.0
% covers: StreamLink, StreamBridge
\datasheet{StreamLink}{The wire under a program that moves 16-byte beats through a wide AXI4 port: each ask one Ethernet frame in, each answer one frame out.}

\dsfacts{\code{lib/parts/src/remote/stream.rs}}{\code{txhdl\_parts::remote::stream::StreamLink}, and \code{StreamBridge}, it and \code{BeatMaster} joined}%
{\code{//docs:examples}}

\subsection*{Function}

\code{HostLink} and \code{LiteMaster} carry 32-bit words to an AXI-Lite
port. FireSim's CPU-managed streams need more: its host driver moves
them with \code{cpu\_managed\_axi4\_read} and
\code{cpu\_managed\_axi4\_write}, on an AXI4 port 128 bits wide with
64-bit addresses (issue 1984). \code{StreamLink} takes a program's ask
for one such beat off the wire and passes it on \code{ask} as a
\code{BeatAsk}; the \code{BeatAnswer} that comes back on \code{back}
leaves as a frame. \code{BeatMaster} performs the ask on the port.

The frames have their own EtherType, \code{0x88b6}, the other one IEEE
leaves for local use, so \code{EthShare} can tell them from
\code{HostLink}'s \code{0x88b5} frames by the type alone.
\code{StreamLink} checks the type itself as well, and ignores a frame
of any other: in \code{rocket\_link}, \code{EthShare} sends it every
frame whose type is not \code{0x88b5}.

\code{StreamBridge} is \code{StreamLink} and \code{BeatMaster} as one
unit, frames on one side and the AXI4 port on the other, joined inside
for the reason \code{HostBridge} is: joined by wires in a hand-written
top, each channel's valid would depend on its own ready through the
other unit. It has no tables of its own here: its ports are
\code{StreamLink}'s \code{rx} and \code{tx} and \code{BeatMaster}'s
port, below and on the next sheet.

\subsection*{Parameters}

\begin{center}\footnotesize
\begin{tabular}{@{}lp{0.7\textwidth}@{}}
\toprule
Parameter & Meaning \\
\midrule
\code{DEV} & The device number, which every ask this accepts must carry
and every answer it sends carries, as for \code{HostLink}. The tables
use 5, as \code{ex\_stream} and \code{//firesim/ax7a200b} do. \\
\bottomrule
\end{tabular}
\end{center}

\subsection*{The frame}

44 bytes, before the padding the sender's MAC adds to reach sixty.
The answer has the same layout, with the differences the table names.

\begin{center}\footnotesize
\begin{tabular}{@{}llp{0.52\textwidth}@{}}
\toprule
Bytes & Field & What \\
\midrule
0 to 13 & header & Ethernet's; type \code{0x88b6}. An answer goes to the
broadcast address from \code{02:00:00:00:00:DEV}. \\
14 & kind & \code{0x03} asks, \code{0x83} answers. \\
15 & device & \code{DEV}. \\
16 & tag & The ask's, which the answer repeats. \\
17 & flags & On an ask, bit 0 a write and bit 1 a no-op; on an answer,
\code{0x01} if the transfer failed and 0 if not. \\
18 to 25 & address & 64 bits, most significant byte first; zero on an
answer. \\
26 to 41 & beat & In memory order: byte 26 is lane 0. On an answer, what
a read read. \\
42 to 43 & strobe & A bit per lane, most significant byte first; zero on
an answer. \\
\bottomrule
\end{tabular}
\end{center}

\dstables{StreamLink}

\subsection*{Behaviour}

\begin{itemize}
\item Every byte coming in is taken, and four are checked as they pass:
the two of the type, the kind and the device. A mismatch sets
\code{alien}, which the last byte clears.
\item \code{got} counts the bytes of the frame coming in and stops at
63 rather than wrapping, so a long frame's later bytes are never read
as its header or its fields again.
\item The address and the strobe are shift registers filled from the
bottom, most significant byte first; the beat is filled from the top,
lane 0 first, so lane 0 ends lowest.
\item An ask goes on \code{ask} at the end of a frame nothing objected
to. A frame that ends by byte 43 is not a whole ask, since the strobe's
last byte is held a cycle after it arrives; the sender's MAC pads every
frame to sixty bytes, so a real ask always goes on past it. With no
room on \code{ask} the ask is dropped, and the program asks again.
\item An answer is taken from \code{back} when no frame is going out,
and its bytes leave one a cycle that \code{tx} has room, spelled from
registers by \code{at}.
\end{itemize}

\subsection*{Verification}

\begin{itemize}
\item Nine tests in \code{lib/parts/src/remote/stream.rs}, in
\code{//lib/parts:parts\_test}, run the two units with a port that keeps
a queue. Among them: \code{a\_wide\_address\_arrives\_whole};
\code{a\_partial\_write\_keeps\_its\_strobe}; and
\code{other\_frames\_are\_ignored}, which sends another device's ask
and another protocol's frame.
\item \code{//lib/examples:ex\_stream} runs a program's asks over a
wire, and the build simulates the three netlists against that run under
nvc and Verilator, among them\newline
\code{//docs:streamlink\_sim\_stream\_link\_stream\_link\_tb\_test}\newline
\code{//docs:streamlink\_sim\_stream\_bridge\_stream\_bridge\_tb\_test}\newline
and their \code{vsim} twins.
\end{itemize}

\subsection*{Used by}

\code{ex\_stream}, and \code{//firesim/ax7a200b}'s \code{rocket\_link},
where \code{EthShare} sends the control frames to a \code{HostBridge}
and every other frame to a \code{StreamBridge<5>}, the stream port of
FireSim's Rocket design on the AX7A200B (issue 1988).

\subsection*{Limits}

One answer frame goes out at a time, and an ask that finds \code{ask}
full is dropped; a program that does not hear an answer asks again.
A frame carries one beat, so a stream moves at a beat per round trip.
Frames carry no authentication.
